\section{Results}
\label{sec:results}

Table~\ref{tab:estimates} gives the fixed-effect estimates. Song onset advanced by 4.1 minutes for each tenfold increase in illuminance, and the interval excludes zero by a wide margin. The noise coefficient was small and its interval spans zero. Dropping noise from the model changed $\beta_L$ by less than 0.1 minutes.

\begin{table}[t]
  \centering
  \caption{Fixed-effect estimates from the model in Equation~\ref{eq:model}. Onset is in minutes relative to civil dawn, so negative values mean earlier song.}
  \label{tab:estimates}
  \begin{tabular}{lrrr}
    \toprule
    Term & Estimate & 95\% CI & $p$ \\
    \midrule
    Illuminance ($\log_{10}$ lux) & $-4.1$ & $[-5.3, -2.9]$ & $<0.001$ \\
    Noise (dB)                    & $-0.04$ & $[-0.13, 0.05]$ & $0.41$ \\
    Temperature ($^\circ$C)       & $-0.9$ & $[-1.1, -0.7]$ & $<0.001$ \\
    Day of year                   & $-0.12$ & $[-0.15, -0.09]$ & $<0.001$ \\
    \bottomrule
  \end{tabular}
\end{table}


Across the observed range of illuminance, from 0.01 to 31 lux, the fitted difference between the darkest and brightest sites is 14.3 minutes. The raw difference is larger, at 27 minutes, because the brightest sites are also warmer.

The piecewise model fits better than the linear one ($\Delta\mathrm{AIC} = 11.6$). The estimated breakpoint is $\hat{\tau} = 3.2$ lux. Below it the slope is 5.6 minutes per tenfold increase, and above it the slope is indistinguishable from zero. Most residential streets in the city are lit to between 5 and 15 lux, which places them on the flat part of the curve.
